\documentclass[11pt]{article}
% Shared preamble for the rigor documents (Lamport-style structured proofs).  \input this file after \documentclass.
\usepackage[T1]{fontenc}\usepackage{lmodern}\usepackage{microtype}
\usepackage[margin=1in]{geometry}
\usepackage{amsmath,amssymb,amsthm,mathtools}
\usepackage{xcolor}\usepackage{enumitem}\usepackage{booktabs}\usepackage{graphicx}\usepackage{hyperref}
\usepackage[most]{tcolorbox}
\definecolor{navy}{HTML}{17365D}\definecolor{teal}{HTML}{117A8B}\definecolor{warn}{HTML}{A33A2B}\definecolor{okgreen}{HTML}{2E7D4F}
\hypersetup{colorlinks=true,linkcolor=navy,citecolor=teal,urlcolor=teal}
\theoremstyle{plain}
\newtheorem{theorem}{Theorem}[section]\newtheorem{lemma}[theorem]{Lemma}\newtheorem{proposition}[theorem]{Proposition}\newtheorem{corollary}[theorem]{Corollary}
\theoremstyle{definition}\newtheorem{definition}[theorem]{Definition}\newtheorem{remark}[theorem]{Remark}\newtheorem{claim}[theorem]{Claim}\newtheorem{issue}[theorem]{Issue}
% ---------- Lamport structured proofs ----------
% \lstep{level}{number}{assertion}   -> indented "<level>number. assertion"
% \lproof                             -> "PROOF:" line (start of the sub-proof of the preceding step)
% \lby{justification}                 -> "BY ..." leaf justification for the preceding step
% \lqed{level}{number}                -> QED step of a sub-proof
% keywords: \lassume \lprove \lsuffices \lcase \ldefine \lpick \lnew
\newlength{\lindent}\setlength{\lindent}{1.7em}
\newcommand{\lstep}[3]{\par\smallskip\noindent\hspace*{\dimexpr\numexpr#1-1\relax\lindent\relax}{\color{navy}$\langle#1\rangle#2$.}\ #3}
\newcommand{\lproof}{\par\noindent\hspace*{\lindent}{\scshape Proof:}}
\newcommand{\lby}[1]{\par\noindent\hspace*{\lindent}{\scshape By}\ #1.}
\newcommand{\lqed}[2]{\par\smallskip\noindent\hspace*{\dimexpr\numexpr#1-1\relax\lindent\relax}{\color{navy}$\langle#1\rangle#2$.}\ {\scshape Q.E.D.}}
\newcommand{\lassume}{{\scshape Assume}\ }\newcommand{\lprove}{{\scshape Prove}\ }\newcommand{\lsuffices}{{\scshape Suffices}\ }
\newcommand{\lcase}{{\scshape Case}\ }\newcommand{\ldefine}{{\scshape Define}\ }\newcommand{\lpick}{{\scshape Pick}\ }\newcommand{\lnew}{{\scshape New}\ }
% ---------- verdict boxes ----------
\newtcolorbox{verdictok}{colback=okgreen!8,colframe=okgreen,title={\bfseries Verdict: verified},fonttitle=\small}
\newtcolorbox{verdictgap}{colback=orange!10,colframe=orange!80!black,title={\bfseries Verdict: gap / caveat},fonttitle=\small}
\newtcolorbox{verdictbreak}{colback=warn!10,colframe=warn,title={\bfseries Verdict: proof-breaking issue},fonttitle=\small}
\newtcolorbox{citedbox}[1]{colback=navy!5,colframe=navy!60,title={\bfseries Cited result: #1},fonttitle=\small,breakable}
\setlength{\parindent}{0pt}\setlength{\parskip}{4pt}\allowdisplaybreaks
\newcommand{\cA}{\mathcal A}\newcommand{\cF}{\mathcal F}\newcommand{\cM}{\mathfrak M}\newcommand{\cN}{\mathfrak N}

% extra calligraphic shortcuts (\providecommand: no clash if a document defines its own)
\providecommand{\cB}{\mathcal B}\providecommand{\cC}{\mathcal C}\providecommand{\cD}{\mathcal D}\providecommand{\cE}{\mathcal E}\providecommand{\cG}{\mathcal G}\providecommand{\cH}{\mathcal H}\providecommand{\cI}{\mathcal I}\providecommand{\cJ}{\mathcal J}\providecommand{\cK}{\mathcal K}\providecommand{\cL}{\mathcal L}\providecommand{\cO}{\mathcal O}\providecommand{\cP}{\mathcal P}\providecommand{\cQ}{\mathcal Q}\providecommand{\cR}{\mathcal R}\providecommand{\cS}{\mathcal S}\providecommand{\cT}{\mathcal T}\providecommand{\cU}{\mathcal U}\providecommand{\cV}{\mathcal V}\providecommand{\cW}{\mathcal W}\providecommand{\cX}{\mathcal X}\providecommand{\cY}{\mathcal Y}\providecommand{\cZ}{\mathcal Z}
\newcommand{\Fid}{\operatorname{F}}\newcommand{\Tr}{\operatorname{Tr}}\newcommand{\eps}{\varepsilon}\newcommand{\dd}{\,\mathrm d}

\providecommand{\one}{\mathbf 1}\providecommand{\norm}[1]{\lVert #1\rVert}
\newcommand{\Nst}[1]{\lVert #1\rVert_{\star}}\newcommand{\Pinch}{\mathcal P}
\newcommand{\Qt}{\widetilde{Q}}\newcommand{\HS}{\mathfrak{S}_2}
\newcommand{\kc}{\kappa_c}\newcommand{\zt}{\zeta}\newcommand{\Phiw}{\Phi_W}

\title{Referee report on \texttt{rigor/certified\_numerics.tex} (agent RC)\\
\large quadratic block bound, cross-term removal, ball-arithmetic enclosures}
\author{Referee RCr (Phase 2b)}
\date{2026-09-09}

\begin{document}
\maketitle

\section{Summary of the verdicts}\label{sec:summary}
% filled last
\providecommand{\eps}{\varepsilon}\providecommand{\Fid}{F}\providecommand{\Tr}{\operatorname{Tr}}
\providecommand{\dd}{\mathrm{d}}\providecommand{\cF}{\mathcal F}
\newcommand{\Xo}{\Xi_0}\newcommand{\Xt}{\Xi_1}

\noindent\textbf{Documents refereed.} \texttt{rigor/certified\_numerics.tex} (agent RC, 2026-09-09,
10~pp.) together with \texttt{rigor/tail\_bound.tex} (agent PT) on which it builds, and the scripts
\texttt{numerics/certified/cert\_qbox.py}, \texttt{cert\_tail\_np.py}, \texttt{cert\_tail.py} and the
data file \texttt{numerics/certified/qball\_Lam60\_k30.pkl}.  All scripts were rerun by the referee
(\S\ref{sec:rerun}).

\begin{center}\footnotesize
\begin{tabular}{p{0.30\textwidth} p{0.11\textwidth} p{0.50\textwidth}}\toprule
Ingredient & Verdict & One-line reason\\\midrule
(a1) $2\times2$ counterexample, Prop.~2.4 & OK & re-derived and reproduced numerically; ratio
 $\to3/\eps$ confirmed to $4$ digits.\\
(a2) Sylvester identity + weight, Thm.~2.2 & OK & identity \eqref{eq:xi} verified symbolically and
 numerically ($10^{-15}$); the derivative formula \eqref{eq:LA} is correct.\\
(a3) Schur bound $\Nst{X_{12}}^2\le\tau$, Cor.~2.6 & GAP & the proof splits the eigenbasis of
 $X_{22}$ into $E^+$ and $E^-$ parts, which presupposes $[X_{22},E^\pm]=0$; and the $E^-$ weight is
 $\le(1-\tau_-)^{-1}$, not $\le1$.  Conclusion is true in all tested cases.\\
(b1) block diagonality of the optimal $U$ & OK & attained by $U_1\oplus U_2$ because $\Fid$
 factorises over the two blocks; the note's one-line justification is thin but repairable.\\
(b2) $\Pinch\sqrt B\le\sqrt{\Pinch B}$ & OK & operator Jensen, right direction.\\
(b3) $\Tr[\sqrt A\,\Pinch\Xi]=-\frac12\|\Xi\|_2^2$ & OK & exact; re-derived and verified numerically
 to $10^{-14}$.  Needs $A=\Pinch(B)$, which holds.\\
(c) $0\le\Phi-\Phi_W\le\Phi_{22}-\log|\det(1+T)|$ & OK & every step in the right direction; the lower
 bound is data processing and is used correctly.\\
(d) enclosures of Table~3 & GAP & the numbers reproduce exactly, but the lower end $\Phi_W$ is
 \emph{not} recomputed from the ball $Q_N$: it is read from \texttt{results\_hp2.txt}, i.e.\ from the
 double-precision trapezoid $Q_N$ that \S5 of the same note declares inaccurate.\\
(e) ball enclosure of $Q_N$ & OK/MINOR & rigorous and reproduced ($5.620\times10^{-18}$,
 $4.687\times10^{-11}$, $91/8281$); the ``$8\,\%$'' reading of the trapezoid error is wrong --- with
 the note's own conversion it is $22\,\%$ at $\zeta=1/15$ and $204\,\%$ at $\zeta=1/120$.\\
(f) statement of what is certified & GAP & the $\widehat D$ dagger is honest; the box cutoff is
 called an ``estimate'' in the verdict but the corresponding column sits inside a table headed
 ``certified enclosures'', and the $Q_N$ contamination of $\Phi_W$ is not disclosed.\\
\bottomrule\end{tabular}
\end{center}

\medskip\noindent\textbf{Proof-breaking issues: none.}  No theorem of
\texttt{certified\_numerics.tex} is false as stated.  Two claims must be weakened: the width of
Table~3 (ingredient (d)) and the ``$8\,\%$ margin'' sentence of \S5 (ingredient (e)).

\section{(a) The square-root bound}\label{sec:a}

\subsection{The counterexample (Prop.~2.4): correct}

\begin{lemma}[referee's re-derivation]\label{lem:cex}
Let $0<\delta<\eps$, $\eps+\delta\le1$, $v=(\sqrt\eps,\sqrt\delta)^{\!\top}$, $B=vv^*$,
$A=\Pinch(B)=\mathrm{diag}(\eps,\delta)$.  Then $0\le B\le1$, $\|X_{12}\|_2^2=\eps\delta$ and
\[
 \frac{\|\sqrt B-\sqrt A\|_2^2}{\|X_{12}\|_2^2}\ \xrightarrow[\delta/\eps\to0]{}\ \frac3\eps .
\]
\end{lemma}
\lstep{1}{1}{$B=vv^*\ge0$, $\|B\|=\|v\|^2=\eps+\delta\le1$, and $\sqrt B=\|v\|^{-1}B$.}
\lby{$B^2=\|v\|^2B$, so $(\|v\|^{-1}B)^2=B$, and $\|v\|^{-1}B\ge0$; uniqueness of the positive square
root}
\lstep{1}{2}{$\|\sqrt B-\sqrt A\|_2^2=\bigl(\tfrac{\eps}{\sqrt{\eps+\delta}}-\sqrt\eps\bigr)^2
 +\bigl(\tfrac{\delta}{\sqrt{\eps+\delta}}-\sqrt\delta\bigr)^2+\tfrac{2\eps\delta}{\eps+\delta}$.}
\lby{$\langle1\rangle1$, $\sqrt A=\mathrm{diag}(\sqrt\eps,\sqrt\delta)$, and
$\|Z\|_2^2=\sum_{ij}|Z_{ij}|^2$; the two off-diagonal entries of $\sqrt B$ are
$\sqrt{\eps\delta}/\sqrt{\eps+\delta}$ and those of $\sqrt A$ vanish}
\lstep{1}{3}{With $r=\delta/\eps\to0$ the three terms are $\frac{\delta^2}{4\eps}(1+O(r))$,
$\delta(1-\sqrt r+O(r))^2$ and $2\delta(1+O(r))$, so the sum is $3\delta(1+O(\sqrt r))$.}
\lby{$\frac{\eps}{\sqrt{\eps+\delta}}-\sqrt\eps=\sqrt\eps\bigl((1+r)^{-1/2}-1\bigr)
=-\frac{\sqrt\eps\,r}2(1+O(r))$ and
$\frac{\delta}{\sqrt{\eps+\delta}}-\sqrt\delta=-\sqrt\delta(1-\sqrt r(1+O(r)))$}
\lqed{1}{4}\lby{$\langle1\rangle3$ divided by $\|X_{12}\|_2^2=\eps\delta$}

\begin{verdictok}
Proposition~2.4 of \texttt{certified\_numerics.tex} is correct.  The referee reproduced the ratio
numerically: $(\eps,\delta)=(0.4,10^{-6})$ gives $7.4921$ against $3/\eps=7.5$, and
$(0.1,10^{-8})$ gives $29.9937$ against $30$.  Since $\eps=\lambda_{\min}(X_{11})=e^{-\kappa_c}$ in
the application, an unweighted bound would carry the factor $e^{\kappa_c}$ and be useless; PT's open
item is therefore genuinely unanswerable in the form in which it was posed, and RC is right to say
so.  Two caveats, both cosmetic: (i) the example is rank one, so $\beta_{\min}(B)=0$ and the
denominators of the exact formula degenerate --- the counterexample says nothing about the
\emph{weighted} statement, as RC's own remark notes; (ii) the hypothesis printed in the proposition
is ``$0<\delta<\eps<1/2$'' whereas $B\le1$ only needs $\eps+\delta\le1$; harmless.
\end{verdictok}


\subsection{The Sylvester identity and the Daleckii--Krein weight (Thm.~2.2): correct}

The identity under review is, for $A,B>0$ with orthonormal eigendecompositions $A\psi_i=a_i\psi_i$,
$B\phi_\alpha=\beta_\alpha\phi_\alpha$ and $V=B-A$,
\begin{equation}\label{eq:xi}
 \bigl\|\sqrt B-\sqrt A\bigr\|_2^2=\sum_{\alpha,i}
 \frac{|\langle\phi_\alpha,V\psi_i\rangle|^2}{(\sqrt{\beta_\alpha}+\sqrt{a_i})^2},
\end{equation}
and the first-order (Daleckii--Krein) formula
\begin{equation}\label{eq:LA}
 \bigl(\mathcal L_A(V)\bigr)_{ij}=\frac{V_{ij}}{\sqrt{a_i}+\sqrt{a_j}},\qquad
 \mathcal L_A(V)=\frac{\dd}{\dd\varepsilon}\Big|_0\sqrt{A+\varepsilon V}.
\end{equation}

\lstep{1}{1}{$\sqrt B\,\Xi+\Xi\sqrt A=B-A$ for $\Xi=\sqrt B-\sqrt A$; hence
$\langle\phi_\alpha,V\psi_i\rangle=(\sqrt{\beta_\alpha}+\sqrt{a_i})\langle\phi_\alpha,\Xi\psi_i\rangle$.}
\lby{$\sqrt B(\sqrt B-\sqrt A)+(\sqrt B-\sqrt A)\sqrt A=B-\sqrt B\sqrt A+\sqrt B\sqrt A-A$; no
commutation is used, which is the point of the Sylvester form; then act with the eigenvectors on the
left and right}
\lstep{1}{2}{\eqref{eq:xi} follows because $(\phi_\alpha)$ and $(\psi_i)$ are complete orthonormal
systems and $\sqrt{\beta_\alpha}+\sqrt{a_i}>0$.}
\lby{$\|\Xi\|_2^2=\sum_{\alpha i}|\langle\phi_\alpha,\Xi\psi_i\rangle|^2$, valid for any two
orthonormal bases; $\langle1\rangle1$}
\lqed{1}{3}\lby{$\langle1\rangle1$, $\langle1\rangle2$}

\begin{verdictok}
Theorem~2.2 of \texttt{certified\_numerics.tex} is correct, and its proof is complete.  The referee
verified \eqref{eq:xi} on four random pairs of $9\times9$ positive matrices (agreement
$\le1.1\times10^{-14}$ in absolute value on quantities of size $8$), and \eqref{eq:LA} together with
$\|\mathcal L_A(V)\|_2^2=2\Nst{W}^2$ on a random block-positive $B$ (agreement $5\times10^{-16}$).
Three points of presentation.  (i) In the summary of \S1 the displayed identity has the numerator
$|\langle\phi_\alpha,X_{12}\psi_i\rangle+\langle\phi_\alpha,X_{12}^*\psi_i\rangle|^2$ and the text
just above it introduces $a_i$ as the eigenvalues of $X_{11}$ and $b_j$ as those of $X_{22}$, while
the formula then uses $a_i$ for the eigenvalues of the \emph{pinching} $A=X_{11}\oplus X_{22}$; the
theorem itself is stated correctly, but the summary should be brought into line, since a reader who
takes $a_i\in\mathrm{spec}(X_{11})$ at face value will read a false statement.  (ii) The identity
requires $A,B>0$ (not merely $\ge0$): with $\beta_\alpha=a_i=0$ and $\langle\phi_\alpha,V\psi_i\rangle
\ne0$ the right-hand side is $+\infty$ while the left-hand side is finite, and the step
$\langle1\rangle1$ then only gives $0=0$.  The theorem says so; the boxed formula does not, and it is
the boxed formula that is quoted in \S1.  (iii) The remark ``no bound of the form
$\|\Xi\|_2\le c\|V\|_2$ can be better than \ldots\ $\lambda_{\min}^{-1/2}$'' is loose (the optimal
constant is $\max_{\alpha i}(\sqrt{\beta_\alpha}+\sqrt{a_i})^{-1}$, attained only if $V$ is supported
on the extremal pair), but the conclusion drawn from it is right.
\end{verdictok}

\subsection{Positivity and the Schur complement (Cor.~2.6): a gap}\label{ss:schur}

Corollary~2.6 asserts $\Nst{X_{12}}^2\le\tau_++\tau_-=\tau$ with PT's occupation tails
$\tau_+=\Tr[E^+\Qt E^+]$, $\tau_-=\Tr[E^-(1-\Qt)E^-]$ of \texttt{tail\_bound.tex}~(1.1); the referee
confirms that the $\tau$ appearing here \emph{is} PT's $\tau$, so the corollary does what the note
claims it does (the weight ``costs nothing'').  The two Schur-complement steps are correct:

\lstep{1}{1}{$X_{12}^*X_{11}^{-1}X_{12}\le X_{22}$ and $X_{12}^*(1-X_{11})^{-1}X_{12}\le1-X_{22}$.}
\lby{the first is the Schur complement of $\Qt\ge0$ with $X_{11}>0$, exactly as in RC's
$\langle1\rangle1$; the second is the same statement for $1-\Qt\ge0$, whose blocks are
$1-X_{11}>0$, $-X_{12}$, $1-X_{22}$}
\lstep{1}{2}{Compressing $\langle1\rangle1$ with $E^+$ resp.\ $E^-$ and tracing gives
\eqref{eq:schur2}.}
\lby{$X_{12}E^+=E\Qt E^+=W^+$, $\Tr[E^+X_{22}E^+]=\tau_+$; $X_{12}E^-=-E(1-\Qt)E^-$,
$\Tr[E^-(1-X_{22})E^-]=\tau_-$; $Z\le Z'\Rightarrow\Tr[P^*ZP]\le\Tr[P^*Z'P]$}
\begin{equation}\label{eq:schur2}
 \Tr\bigl[(W^+)^*X_{11}^{-1}W^+\bigr]\le\tau_+,\qquad
 \Tr\bigl[(W^-)^*(1-X_{11})^{-1}W^-\bigr]\le\tau_- .
\end{equation}

\noindent The gap is in step $\langle1\rangle3$ of the note, which reads
\begin{quote}\itshape
$\Nst{X_{12}}^2=\sum_{i,j\in E^+}\frac{|V_{ij}|^2}{(\sqrt{a_i}+\sqrt{b_j})^2}
+\sum_{i,j\in E^-}\frac{|V_{ij}|^2}{(\sqrt{a_i}+\sqrt{b_j})^2}$,
\end{quote}
justified by ``on the $E^+$ part drop $\sqrt{b_j}$ \ldots\ on the $E^-$ part $b_j\ge1-\tau_-\ge1/2$''.

\begin{verdictgap}
Two defects, both repairable, in the proof of Corollary~2.6.  (i) The index $j$ runs over an
eigenbasis of $X_{22}=E^\perp\Qt E^\perp$, and $i$ over an eigenbasis of $X_{11}$; the displayed
splitting of the sum into an ``$E^+$ part'' and an ``$E^-$ part'' therefore presupposes that
$X_{22}$ commutes with $E^+$ and $E^-$, i.e.\ that $E^+\Qt E^-=0$, which is nowhere assumed and is
false in general ($E^\pm$ are spectral projections of the \emph{vacuum} modular Hamiltonian $K$, not
of $\Qt$).  The same presupposition underlies ``the eigenvalues of $X_{22}$ on $E^-H$ are within
$\tau_-$ of $1$''.  (ii) Even granting the splitting, on the $E^-$ part
$(\sqrt{a_i}+\sqrt{b_j})^{-2}\le b_j^{-1}\le(1-\tau_-)^{-1}$, not $\le1$; with the crude bound
$b_j\ge1/2$ actually quoted one gets $\le2$.  A clean repair that avoids both defects: split the
eigenbasis of $X_{22}$ by $b_j<\frac12$ / $b_j\ge\frac12$ (spectral projections $R_<,R_\ge$ of
$X_{22}$, hence commuting with $X_{22}$ by construction), use the weight $a_i^{-1}$ and
\eqref{eq:schur2} on $R_<$, and the weight $2$ together with
$\|X_{12}R_\ge\|_2^2=\|E(1-\Qt)E^\perp R_\ge\|_2^2\le\|1-X_{11}\|\Tr[R_\ge(1-X_{22})R_\ge]$ on
$R_\ge$; this yields $\Nst{X_{12}}^2\le\Tr[R_<X_{22}R_<]+2\Tr[R_\ge(1-X_{22})R_\ge]$, at the price of
a factor $2$ and of a small additional argument comparing that quantity to $\tau$.  The conclusion
itself survives every test: on the free-fermion data ($\zeta=1/15$, $\kappa_c=25.3$,
$\Lambda=60$, $\kappa_{\max}=30$) the referee measures $\Nst{X_{12}}^2=1.893\times10^{-7}$ against
$\tau=2.648\times10^{-7}$, ratio $0.71$, and $\|E^+X_{22}E^-\|_2/\|X_{22}\|_2=1.3\times10^{-6}$, so
the omitted coupling is numerically irrelevant here; over $400$ random $12\times12$ positive
contractions with a genuine $E^+/E^-$ mixing the ratio never exceeded $0.15$.
Finally, the \S1 sentence ``\emph{positivity alone} gives $\Nst{X_{12}}^2\le\tau$'' is inaccurate:
the $E^-$ half uses $\Qt\le1$ as well, i.e.\ that $\Qt$ is a positive \emph{contraction}, and the
weight estimate uses the smallness of $\tau_-$.
\end{verdictgap}

\section{(b) Removal of the Bures cross term}\label{sec:b}

\subsection{Is the optimal Uhlmann unitary of the pinched pair block diagonal?}

The claim (step $\langle1\rangle3$ of Thm.~4.2) is: $Q$ commutes with $E$ and $A=\Pinch(\Qt)$ is
block diagonal, hence the maximiser $U$ of $|\det(S_QS_A+C_QUC_A)|$ may be taken block diagonal, and
then $-\log|\det M_A|=\Phi_W+\Phi_{22}$.  The note's justification is one clause
(``$Q$ and $A$ are both block diagonal, so $U$ may be chosen block diagonal'').  It is correct, but
the argument that makes it correct is the \emph{second} half of the same step, and should be stated
in this order:

\lstep{1}{1}{$\omega_Q$ and $\omega_A$ both factorise as $\omega_{Q_{11}}\otimes\omega_{Q_{22}}$ resp.\
$\omega_{X_{11}}\otimes\omega_{X_{22}}$ under $\cF(H)\cong\cF(EH)\otimes\cF(E^\perp H)$.}
\lby{a gauge-invariant quasi-free state with a symbol that is block diagonal for $H=EH\oplus E^\perp H$
is the product of the two quasi-free states with the diagonal blocks (Note~5, proof of Thm.~3.1);
$[Q,E]=0$ and $A=\Pinch(\Qt)$}
\lstep{1}{2}{$\Fid(\omega_Q,\omega_A)=\Fid(\omega_{Q_{11}},\omega_{X_{11}})\cdot
\Fid(\omega_{Q_{22}},\omega_{X_{22}})=e^{-\Phi_W}e^{-\Phi_{22}}$.}
\lby{multiplicativity of the fidelity over tensor products, $\langle1\rangle1$, $EAE=X_{11}=E\Qt E$,
and PT's definition $\Phi_W=-\log\Fid(\omega_{EQE},\omega_{E\Qt E})$}
\lstep{1}{3}{Let $U_1,U_2$ be optimal unitaries for the two blocks in the sense of Note~5, Thm.~2.1.
Then $U:=U_1\oplus U_2$ satisfies $|\det(S_QS_A+C_QUC_A)|=e^{-\Phi_W-\Phi_{22}}=\Fid(\omega_Q,\omega_A)$,
so $U$ is optimal and $M_A$ is block diagonal and invertible.}
\lby{for block-diagonal $U$ the matrix $M_A$ is block diagonal and its determinant is the product of
the two block determinants, each equal to the corresponding block fidelity by Note~5, Thm.~2.1;
$\langle1\rangle2$; and no unitary can exceed $\Fid$ by Thm.~4.2$\langle1\rangle1$}
\lqed{1}{4}\lby{$\langle1\rangle1$--$\langle1\rangle3$}

\begin{verdictok}
Correct, with the caveat that the note's phrasing invites the reader to look for a
symmetry/uniqueness argument for the maximiser, whereas the real argument is that a block-diagonal
$U$ \emph{attains} the (block-multiplicative) maximum, so that one may take it.  Note that
block diagonality of $M_A$ is load bearing: it is what makes $Y_0,Y_1$ block diagonal in
$\langle1\rangle5$ and hence what allows Lemma~3.1 to be applied at all.  If the maximiser were not
block diagonal, the whole mechanism would collapse.  The referee checked the identity
$-\log|\det M_A|=\Phi_W+\Phi_{22}$ numerically on the free-fermion data with $U$ taken from the polar
decomposition of $G_Q^{1/2}G_A^{1/2}$, exactly as \texttt{cert\_tail\_np.py} does
(\S\ref{sec:rerun}).  A second caveat: ``$U_1,U_2$ optimal'' is Note~5, Thm.~2.1, which is quoted in
\texttt{tail\_bound.tex} only through the remark after Lemma~3.1; \texttt{certified\_numerics.tex}
should carry the exact locator, because the entire cross-term removal rests on the
polar-decomposition characterisation of the maximiser.
\end{verdictok}

\subsection{The operator Jensen step, and the trace identity}

\begin{lemma}[referee's check of Lemma 3.1]\label{lem:jens}
Let $B\ge0$ on $\mathfrak h_1\oplus\mathfrak h_2$, $A=\Pinch(B)$, $\Xo=\sqrt B-\sqrt A$,
$\Theta_0=-\Pinch(\Xo)$.  Then $\Theta_0\ge0$ and $\Tr[\sqrt A\,\Theta_0]=\frac12\|\Xo\|_2^2$,
equivalently $\Tr[\sqrt A\,\Pinch\Xo]=-\frac12\|\Xo\|_2^2$.
\end{lemma}
\lstep{1}{1}{$\Theta_0=\sqrt{\Pinch B}-\Pinch\sqrt B\ge0$.}
\lby{$\sqrt A$ is block diagonal so $\Pinch\sqrt A=\sqrt A$; and $\Pinch(f(B))\le f(\Pinch B)$ for
operator concave $f$ and a unital positive map $\Pinch$ (Davis 1957; Hansen--Pedersen 2003,
Thm.~2.1), with $f=\sqrt\cdot$ operator concave on $[0,\infty)$.  The direction is the one used:
concavity gives $\Pinch f(B)\le f(\Pinch B)$, i.e.\ the pinching of the square root is \emph{below}
the square root of the pinching}
\lstep{1}{2}{$\Tr[\sqrt A\,\Pinch\sqrt B]=\Tr[\sqrt A\sqrt B]$ and $\Tr A=\Tr B$.}
\lby{$\Pinch$ is self-adjoint for the Hilbert--Schmidt pairing and $\Pinch\sqrt A=\sqrt A$;
$\Pinch$ is trace preserving}
\lstep{1}{3}{$\|\Xo\|_2^2=\Tr B+\Tr A-2\Tr[\sqrt A\sqrt B]=2\bigl(\Tr A-\Tr[\sqrt A\sqrt B]\bigr)
=2\Tr[\sqrt A\,\Theta_0]$.}
\lby{expand $\Tr[(\sqrt B-\sqrt A)^2]$; $\langle1\rangle2$}
\lqed{1}{4}\lby{$\langle1\rangle1$, $\langle1\rangle3$}

\begin{verdictok}
Lemma~3.1 of \texttt{certified\_numerics.tex} is correct, including \eqref{eq:Ytr} below, and its
proof is complete.  Note that the trace identity is \emph{not} an inequality and uses $A=\Pinch(B)$
twice (through $\Tr A=\Tr B$ and $\Pinch\sqrt A=\sqrt A$); it would fail for a general block-diagonal
$A$.  Numerical check by the referee on three random $9\times9$ positive $B$ with a $4|5$ split:
$\Tr[\sqrt A\Pinch\Xo]$ and $-\frac12\|\Xo\|_2^2$ agree to $1.7\times10^{-14}$ on quantities of size
$1.6$, and $\lambda_{\min}(\Theta_0)>0$ in all three cases.  The consequence used later,
\begin{equation}\label{eq:Ytr}
 Y=Y^*\ \text{block diagonal},\ -c\sqrt A\le Y\le c\sqrt A\ \Longrightarrow\
 |\Tr[Y\Xo]|\le\tfrac c2\|\Xo\|_2^2,
\end{equation}
follows because $\Tr[Y\Xo]=\Tr[Y\Pinch\Xo]=-\Tr[Y\Theta_0]$ and $\Theta_0\ge0$; this is the whole
mechanism: \emph{a block-diagonal observable sees the off-diagonal perturbation only at second
order}.  It is a genuinely new and, in the referee's judgement, correct idea, and it is the reason
the $2\sqrt{\Phi\Phi_\partial}$ of \texttt{tail\_bound.tex} disappears.
\end{verdictok}


\section{(c) The final inequality}\label{sec:c}

The computable statement, \eqref{eq:main} of the note, is
\begin{equation}\label{eq:main}
 0\ \le\ \Phi-\Phi_W\ \le\ \Phi_{22}-\log\bigl|\det(1+T)\bigr|,\qquad
 T=M_A^{-1}\bigl(S_Q\Xo+C_QU\Xt\bigr),
\end{equation}
with $A=\Pinch(\Qt)=X_{11}\oplus X_{22}$, $S_Q=Q^{1/2}$, $C_Q=(1-Q)^{1/2}$, $S_A=A^{1/2}$,
$C_A=(1-A)^{1/2}$, $\Xo=\Qt^{1/2}-S_A$, $\Xt=(1-\Qt)^{1/2}-C_A$, $M_A=S_QS_A+C_QUC_A$ with $U$ the
optimal block-diagonal Uhlmann unitary of the pair $(Q,A)$, and
$\Phi_{22}=-\log\Fid(\omega_{Q_{22}},\omega_{X_{22}})$, $Q_{22}=E^\perp QE^\perp$,
$X_{22}=E^\perp\Qt E^\perp$.  All of these are unambiguous, and match the implementation in
\texttt{cert\_tail\_np.py} line by line.  The chain of inequalities is

\lstep{1}{1}{$\Fid(\omega_Q,\omega_{\Qt})\ \ge\ |\det(S_QS+C_QUC)|$ for the chosen $U$
($\ge$: an arbitrary purification of $\omega_{\Qt}$ is used).}
\lby{Thm.~4.2$\langle1\rangle1$; the isometry $\Theta^U=\binom{S}{UC}$ has
$\Theta^{U*}\Theta^U=1$ and $(\Theta^U\Theta^{U*})_{11}=\Qt$, and monotonicity of the fidelity under
restriction \emph{increases} it, so any purification pair gives a lower bound}
\lstep{1}{2}{$S_QS+C_QUC=M_A(1+T)$ exactly.}
\lby{$S=S_A+\Xo$, $C=C_A+\Xt$, and $M_A$ invertible}
\lstep{1}{3}{$\Phi\le-\log|\det M_A|-\log|\det(1+T)|=\Phi_W+\Phi_{22}-\log|\det(1+T)|$.}
\lby{$\langle1\rangle1$, $\langle1\rangle2$, monotonicity of $-\log$, and \S\ref{sec:b}}
\lstep{1}{4}{$\Phi\ge\Phi_W$.}
\lby{$\Fid$ is monotone under the restriction to the subalgebra $\cF(EH)$, and the restriction of a
quasi-free state is quasi-free with the compressed symbol; $EQE$ and $E\Qt E$ are those symbols
(\texttt{tail\_bound.tex}, Thm.~5.1$\langle1\rangle1$).  This is data processing and is used in the
correct direction: it is the \emph{lower} end of the bracket}
\lqed{1}{5}\lby{$\langle1\rangle3$, $\langle1\rangle4$}

\begin{verdictok}
\eqref{eq:main} is correct as stated, every step points the right way, and the lower bound
$\Phi\ge\Phi_W$ is the standard data-processing inequality used correctly.  Three observations.
(i) \eqref{eq:main} holds with \emph{no} hypothesis --- neither (H) nor $\|T\|<1$ ---
because it never expands the logarithm; the note says this and is right.  Consequently the sign of
$\Phi_{\rm off}=-\log|\det(1+T)|$ is not controlled a priori; it happens to be positive in all seven
data points, but if it were negative the bracket would simply be tighter.  (ii) The a priori form
(H)--(4.4) is also correct: $|\Re\Tr T|\le\frac{c_0}2\|\Xo\|_2^2+\frac{c_1}2\|\Xt\|_2^2$ by
\eqref{eq:Ytr}, and $\log|\det(1+T)|\ge\Re\Tr T-\|T\|_2^2/(2(1-\|T\|))$ by
$|\Tr T^k|\le\|T\|_2^2\|T\|^{k-2}$; the referee checked both directions of the second estimate and the
Hermiticity and block diagonality of $Y_0,Y_1$.  (iii) The remaining weakness is honestly located by
the note: the measured $c_0$ (referee's rerun: $18.4$ to $31.3$ over the seven $\zeta$ at
$\kappa_c=25.3$) comes from the window edge, where $\lambda_{\min}(S_A)$ is three orders below
$\lambda_{\min}(S_Q)$, so (H) cannot hold with $c_0=O(1)$; the a priori theorem is
therefore a factor $4$--$6$, not $50$--$150$, better than \texttt{tail\_bound.tex}.  A remark the
note should add: since $Y_0=\frac12(M_A^{-1}S_Q+S_QM_A^{-*})$ and $M_A\approx S_QS_A+C_QC_A$, the
natural comparison operator for $Y_0$ is not $S_A$ but $S_A+O(\ldots)$; a hypothesis of the form
$\pm Y_0\le c_0'(S_A+\lambda)$ with the window edge $\lambda=\sqrt\eps$ excised would presumably give
$c_0'=O(1)$ and an unconditional $O(\tau)$ theorem.  As it stands, ``$c_0\sim18$--$42$'' is a
placeholder, not a constant, and \eqref{eq:main} is the only statement that should be quoted.
\end{verdictok}

\section{(d) Reproduction of the enclosures, and the box cutoff}\label{sec:rerun}

\paragraph{Reruns.}  All three scripts were rerun by the referee with the venv interpreter.
\texttt{cert\_qbox.py 60 30 200} finishes in $4.4$\,s and prints
\begin{center}\small\ttfamily
Lam=60.0 kmax=30.0 N=91 prec=200  max entry radius = 5.620e-18\\
\ \ max |Q\_ball - Q\_float| = 4.687e-11;  float entries inside ball: 91/8281
\end{center}
reproducing the two numbers of \S5 of the note ($5.6\times10^{-18}$, $4.7\times10^{-11}$) and the
claim that exactly the $N=91$ diagonal float entries lie inside their enclosures.
\texttt{cert\_tail\_np.py} was rerun for all seven $\zeta$ at $\eps=10^{-11}$ ($\kappa_c=25.33$) with
$\Phi_W$ taken, as the note does, from \texttt{results\_hp2.txt}; every entry of the
$\kappa_c=25.3$ half of the table of \S7 is reproduced to the printed digits
($\Phi_{22}=1.936(-9),8.002(-9),3.255(-8),1.313(-7),5.271(-7),2.109(-6),8.383(-6)$;
$\Phi_{\rm off}=2.598(-9),1.905(-8),2.462(-8),1.315(-7),4.938(-7),1.675(-6),6.803(-6)$; widths
$0.79,1.19,0.64,0.76,0.78,0.81,1.00\,\%$), and so is every interval of Table~3, e.g.\
$0.057444+0.0000194+0.0000260=0.057897$ at $\zeta=1/120$ in units $10^{-5}$.  The
containment claim is also reproduced: each ``Note~5/PT best'' value lies inside the last-column
interval (the tightest case is $\zeta=4/15$: $47.2564\in[46.51575,47.42576]$).
\texttt{cert\_tail.py}, the \texttt{mpmath} version that recomputes $\Phi_W$ itself, is present but
its only log file, \texttt{numerics/certified/tail\_z0.2.log}, is \emph{empty}: the high-precision
run was evidently never completed.

\paragraph{The box-cutoff column.}  The last two columns of Table~3 add
$0.0641\,\zeta^2\kappa_{\max,\rm eff}^{-2}$ with $\kappa_{\max,\rm eff}=29.6$; the referee confirms
the arithmetic ($7.316\times10^{-5}\zeta^2$, e.g.\ $5.08\times10^{-9}$ at $\zeta=1/120$, giving
$0.058405$).  The provenance of the constant is \texttt{tail\_bound.tex}~(7.2),
$f_2-f_2^{(\kappa_{\max})}=0.513/(8\kappa_{\max}^2)=0.0641/\kappa_{\max}^2$, which is the
\emph{leading asymptotics} of the $\zeta^2$ coefficient, obtained from the exact leading ultraviolet
occupation tail --- not a proved bound on $\Phi_E-\Phi_{P_{\rm sub}}$.  PT is explicit about this
(\S7b: ``Theorem 4.1 \ldots\ does \emph{not} give a bound on $\Phi_E-\Phi_{P_{\rm sub}}$.
Empirically it is small'').


\subsection{The arithmetic behind Table~3 does not resolve $\Phi_{\rm off}$}\label{ss:noise}

\texttt{cert\_tail\_np.py} evaluates $\Xo,\Xt,T$ and $\log\det(1+T)$ in \emph{double precision}
(the note says so) on matrices whose spectrum spans $[2.8\times10^{-13},1-2.8\times10^{-13}]$.  The
script never checks the one identity that its own theory provides as an error monitor, namely
$-\log|\det M_A|=\Phi_W+\Phi_{22}$ (step $\langle1\rangle3$ of Thm.~4.2).  The referee did:
\begin{center}\small
\begin{tabular}{l r r r}\toprule
input ($\zeta=1/120$, $\kappa_c=25.3$) & $-\log|\det M_A|$ & $\Phi_W+\Phi_{22}$ & $\Phi_{\rm off}$\\\midrule
ball $Q_N$ (midpoints)              & $1.0285\times10^{-4}$ & $5.7638\times10^{-7}$ & $+2.598\times10^{-9}$\\
trapezoid $Q_N$                     & $4.6649\times10^{-5}$ & $5.7638\times10^{-7}$ & $-1.158\times10^{-10}$\\
ball $+$ $10^{-16}$ jitter (3 seeds)& $1.09$--$1.76\times10^{-4}$ & --- & $1.49,2.38,2.94\times10^{-9}$\\
\bottomrule\end{tabular}
\end{center}
The identity is \emph{true}: on the $10\times10$ complement block, where the same construction can be
carried out in \texttt{mpmath}, the referee finds $-\log|\det M_A^{22}|=1.9359858\times10^{-9}$ at
\texttt{dps}$=60$ against $\Phi_{22}=1.935986\times10^{-9}$ --- agreement to $7$ digits, which also
confirms independently that the polar unitary of $G_Q^{1/2}G_A^{1/2}$ is the maximiser and that
\S\ref{sec:b} is right.  In double precision the same $10\times10$ quantity comes out as
$1.98\times10^{-9}$ ($2.3\,\%$ off), and on the full $91\times91$ problem the double-precision value
of $-\log|\det M_A|$ misses $\Phi_W+\Phi_{22}$ by $1.0\times10^{-4}$, i.e.\ by $180$ times the
quantity itself.

\begin{verdictbreak}
\emph{The $\Phi_{\rm off}$ column of the table in \S7, and hence the upper end of every interval of
Table~3, is not supported by the computation that produced it.}  Under a relative jitter of
$10^{-16}$ in the input $Q_N$ --- i.e.\ within the rounding noise of the double-precision pipeline
itself --- the computed $\Phi_{\rm off}$ at $\zeta=1/120$ takes the values
$2.60,\,2.94,\,1.49,\,2.38\times10^{-9}$ and, with the trapezoid input, $-1.16\times10^{-10}$: the
tabulated value $2.60(-9)$ has no significant digit and does not even have a stable sign.  At
$\zeta=1/15$ the spread is $1.05$--$1.75\times10^{-7}$ around the tabulated $1.32\times10^{-7}$,
i.e.\ $\pm25\,\%$.  The reason is the cancellation in $\Re\Tr T$: $\|T\|_2^2$ is perfectly stable
($3.917\times10^{-9}\pm10^{-12}$ under the same jitter), but $\log|\det(1+T)|$ is a difference of
quantities of size $10^{-4}$ in this arithmetic.  Consequently the printed widths
``$0.64$--$1.19\,\%$'' are not established, and \emph{no} interval of Table~3 may be quoted as a
certified upper bound in its present form.  This is not an error in Thm.~4.2 or in
\eqref{eq:main} --- both are correct --- and the fix is already in the repository: the
\texttt{mpmath} script \texttt{cert\_tail.py} recomputes exactly these quantities (and $\Phi_W$
itself) at \texttt{dps}$=40$ from the ball $Q_N$, but its only log,
\texttt{numerics/certified/tail\_z0.2.log}, is empty, so it appears never to have been run to
completion.  The three quantities that \emph{are} stable under the jitter test --- $\Phi_{22}$
(computed at \texttt{dps}$=60$), $\|\Xo\|_2^2,\|\Xt\|_2^2,\|T\|_2^2,\|T\|$ and $c_0,c_1$ (norms, no
cancellation) --- support the \emph{a priori} bound (4.4), which the referee's rerun gives as
$17.8$--$19.9\,\%$ of $\Phi_W$ across the seven $\zeta$ (the note says $20$--$30\,\%$).
\end{verdictbreak}

\begin{verdictgap}
Independently of \S\ref{ss:noise}: the lower end of every interval of Table~3 is
$\Phi_W$ read from \texttt{results\_hp2.txt}, which was produced by \texttt{fidelity\_hp.py} from
\texttt{compression\_box.Q\_box}, i.e.\ from the $4\times10^5$-node double-precision \emph{trapezoid}
$Q_N$ that \S5 of the same note declares inaccurate --- contradicting the sentence ``All numbers
below use the ball $Q_N$''.  The referee recomputed $\Phi_W$ at \texttt{dps}$=60$ from both inputs at
$\zeta=1/120$: $5.744435427\times10^{-7}$ (ball) against $5.744385761\times10^{-7}$ (trapezoid,
reproducing \texttt{results\_hp2.txt} to $8$ digits and thereby confirming the provenance).  The
shift is $+4.94\times10^{-12}$, i.e.\ $8.6\times10^{-6}$ relative --- harmless, and in the safe
direction (the quoted lower end remains a lower bound) --- but it is measured here, not bounded in
the note, and the note's own error formula would have predicted something $10^5$ times larger
(next subsection).
\end{verdictgap}

\section{(e) The ball enclosure of $Q_N$}\label{sec:e}

\paragraph{Rigour of the enclosure.}  The construction is sound.  $Q^{\rm c}_{jj}$ and
$Q^{\rm c}_{jl}$ depend on the $2N$ numbers $I(k_j)$, $C(k_j)$ only, each a one-dimensional integral
of a function that is meromorphic in $t$ with poles exactly at $t\in2\pi i\mathbb Z$, so it is
analytic in $|\operatorname{Im}t|<2\pi$ as claimed and \texttt{acb.integral} (Petras' algorithm with
Gauss--Legendre nodes and a Bernstein-ellipse remainder) returns a rigorous ball on $[d,\Lambda]$.
The Python integrand ignores the \texttt{analytic} flag; this is safe \emph{here} because the only
singularity is a pole and \texttt{arb} division by a ball containing $0$ returns an indeterminate
ball, which forces subdivision --- it would not be safe for a branch cut, and the script deserves a
comment saying so.  The endpoint pieces are enclosed by
$\frac{\sin kt}{\sinh(t/2)}=2k(1+e_1)$, $\frac{\cos kt-1}{\sinh(t/2)}=-k^2t(1+e_2)$ with
$|e_1|\le\frac{(kd)^2}6+\frac{d^2}{24}$, $|e_2|\le\frac{(kd)^2}{12}+\frac{d^2}{24}$; the referee
confirms $\bigl|\frac{\sin x}x-1\bigr|\le\frac{x^2}6$, $\bigl|\frac{x}{\sinh x}-1\bigr|\le\frac{x^2}6$
and $\bigl|\frac{1-\cos x}{x^2/2}-1\bigr|\le\frac{x^2}{12}$ from the alternating Taylor series with
decreasing terms, and that the integrands $2k(\Lambda-t)$ and $-k^2t$ do not change sign on $[0,d]$,
so the products may be enclosed as printed.  Two small slips: the product of the two Taylor factors
contributes $|e_1e_2|$ as well, so strictly $|e|\le|e_1|+|e_2|+|e_1||e_2|$ (the omitted term is
$\sim10^{-25}$ and cannot affect a radius of $5.6\times10^{-18}$); and the claimed scaling ``it
improves as $d^2$'' is wrong --- the endpoint contribution to $I$ is
$\sim2k\Lambda d\cdot(kd)^2/6=O(d^3)$ and to $C$ is $O(d^4)$.  Both err on the safe side.  The
referee reran \texttt{cert\_qbox.py 60 30 200}: $4.4$\,s, max radius $5.620\times10^{-18}$,
$\max|Q^{\rm ball}-Q^{\rm trap}|=4.687\times10^{-11}$, $91/8281$ float entries inside their balls.
With $k_{\max}=2\pi\cdot45/60$ and $d=10^{-6}$ the dominant term is
$2k\Lambda d\cdot(kd)^2/6/(2\pi\Lambda)\approx5.6\times10^{-18}$, so the note's attribution of the
radius to the Taylor cut is right.


\paragraph{The ``$8\,\%$ rather than $3\,\%$'' sentence.}  \S5 of the note reads: ``\emph{it
corresponds to $\eta\approx4\times10^{-8}$ in trace norm, i.e.\ an arithmetic margin of $8\,\%$
rather than the assumed $3\,\%$.}''  The referee checks the chain.  (i)
$\|Z\|_1\le N^{3/2}\max_{jl}|Z_{jl}|$ is correct and gives
$\eta=868.25\times4.687\times10^{-11}=4.07\times10^{-8}$, as claimed.  (ii) Inserting this into PT's
(8.1) ($|\Phi_{\rm sub}-\Phi^{\rm c}_{\rm sub}|\le2\sqrt{2\Phi}\sqrt{5\eta}+5\eta$) gives
$7.5\times10^{-6}$ at $\zeta=1/15$, which is $21.7\,\%$ of $\Phi$, not $8\,\%$; the value $8\,\%$
corresponds to $\eta\approx5.5\times10^{-9}\approx N\max_{jl}|Z_{jl}|$, i.e.\ to a conversion
different from the one the note states, weaker by $\sqrt N\approx9.5$.  (iii) More importantly the
margin PT~(8.1) scales like $\sqrt\Phi$, so the \emph{relative} margin scales like
$\Phi^{-1/2}\propto\zeta^{-1}$ and no single percentage describes it: with the note's own $\eta$ the
margin is $3.3\,\%$ at $\zeta=8/15$, $21.7\,\%$ at $\zeta=1/15$ and $204\,\%$ at $\zeta=1/120$.
The answer to the question ``does an entry error of $4.7\times10^{-11}$ translate into $8\,\%$ of
$\Phi$ at the smallest $\zeta$'' is therefore \emph{no}: by the note's own formula it is a factor
$25$ worse there.  (iv) That bound is, however, wildly pessimistic: the referee's direct
recomputation of $\Phi_W$ from the two inputs (\S\ref{ss:noise}) gives a relative shift of
$8.6\times10^{-6}$ at $\zeta=1/120$, i.e.\ $2\times10^{5}$ times smaller than
PT~(8.1) allows, because Powers--St\o rmer discards the fact that the perturbation of $Q_N$
is smooth and almost commutes with the state.  The correct summary of \S5 is: \emph{the trapezoid
$Q_N$ carries an entrywise error of $4.7\times10^{-11}$, for which the only available rigorous
sensitivity estimate is useless at small $\zeta$ ($3$--$204\,\%$); the ball $Q_N$ removes the
question entirely, with a margin $\le5.9\times10^{-4}$ relative at every tabulated $\zeta$.}
The claim that the ball reduces PT's margin ``by a factor $400$'' is correct at $\zeta=1/15$
($\eta_Q\le4.9\times10^{-15}$, margin $2.6\times10^{-9}$).

\begin{verdictok}
The ball-arithmetic enclosure of $Q_N$ (\S5 of the note) is rigorous, reproduces exactly, and is the
strongest part of Part~II.  The two defects are editorial: the ``$8\,\%$'' figure is not what the
note's own conversion and PT's (8.1) give (it should read $3$--$204\,\%$ across the
tabulated $\zeta$, or, honestly, ``the rigorous sensitivity bound is vacuous at small $\zeta$ and the
measured shift is $10^{-5}$ relative''), and the endpoint error improves as $d^3$, not $d^2$.
\end{verdictok}

\small
\section{(f) Is the statement of what is certified honest and complete?}\label{sec:f}

\begin{verdictgap}
Partly.  Honest and clearly flagged: the $\widehat D$ quadrature (dagger on Table~3, \S6, with two
concrete routes to a proof and the empirical $10^{-13}$ convergence test), the constant
$c_0\in[18,42]$ and the resulting $4$--$6\times$ weakening of the a priori theorem, and the
sentence naming the box cutoff as the dominant remaining uncertainty.  Not flagged, and needed
before the note is used:
\begin{enumerate}[label=(\roman*),leftmargin=1.8em]
\item \emph{The arithmetic of Table~3 itself.}  \S\ref{ss:noise}: $\Phi_{\rm off}$ is computed in
 double precision and is noise at the smallest $\zeta$; the note states the precision but draws no
 consequence, and no error term for it appears in the table.  A ``certified enclosure'' whose upper
 end is a float64 log-determinant of a matrix with condition number $\sim10^{13}$ needs either
 \texttt{cert\_tail.py} (present, apparently never completed) or an interval evaluation.
\item \emph{The provenance of $\Phi_W$.}  ``All numbers below use the ball $Q_N$'' is false for the
 lower end of Table~3, which comes from \texttt{results\_hp2.txt} (trapezoid).  The effect is
 $8.6\times10^{-6}$ relative and in the safe direction, but this must be said, not left to the
 reader to discover.
\item \emph{The ball radius is never propagated.}  $5.6\times10^{-18}$ per entry is negligible, but
 no step of \S7 carries it, so ``certified'' is used for a computation with no interval arithmetic
 in it after \S5.
\item \emph{Overclaims in the scope section.}  \S1 announces (a) an unconditional bound
 $\Phi-\Phi_W\le C_\star(\tau+\nu)/(1-C_\star(\tau+\nu))$ with $C_\star=6$, hence
 $\Phi-\Phi_W\le0.85\,\zeta^2\kappa_c^{-2}$; no such statement occurs in the body, and the body's a
 priori bound carries $c_0\approx35$ instead.  (b) ``the spectral window by verified eigenvalue
 enclosures (Rayleigh--Ritz $+$ residual, \S6)'' --- \S6 proposes Gershgorin discs on $U^*Q_NU$ and an
 a posteriori Sylvester bound, and \texttt{cert\_tail\_np.py} uses neither (it calls
 \texttt{numpy.linalg.eigh}).  (c) ``the fidelity determinant of Note~5, Thm.~2.1 in ball
 arithmetic'' --- not done anywhere.  (d) ``the \emph{measured} (and, in Part~II, certified) values
 of $\tau$'' --- $\tau$ is not certified in Part~II.
\item \emph{The box cutoff.}  The verdict calls it an estimate, but the column ``$\Phi$ incl.\ box
 cutoff'' sits in a table captioned ``Certified enclosures'' and the verdict itself says
 ``the certified brackets \ldots\ are $1.7$--$2.4\,\%$ wide for $\Phi$ itself once PT's box-cutoff
 estimate is added''.  Adding an unproved (asymptotic, one-term) correction to an interval does not
 produce a certified interval; that column should be labelled ``extrapolated''.  (Its width range is
 also $1.5$--$2.4\,\%$, not $1.7$--$2.4\,\%$: the $\zeta=1/30$ row is $1.54\,\%$.)
\end{enumerate}
\end{verdictgap}

\section{What may be quoted, and what must not}\label{sec:quotable}

\begin{enumerate}[leftmargin=1.8em]
\item \textbf{Quotable as proved (analysis).}  Theorem~2.2 (the Sylvester identity and the
 Daleckii--Krein weight), Proposition~2.4 (the unweighted bound is false, with the sharp rate
 $\eps^{-1}$), Lemma~3.1 (operator Jensen $+$ the trace identity), Theorem~4.2 and its
 hypothesis-free corollary \eqref{eq:main}.  Corollary~2.6 ($\Nst{X_{12}}^2\le\tau$) needs the
 repair of \S\ref{ss:schur} before being quoted with a constant $1$; it is safe with a factor $2$.
\item \textbf{Quotable as certified (arithmetic).}  The enclosure of the box symbol:
 $|Q^{\rm ball}_{jl}-Q_{jl}|\le5.62\times10^{-18}$ at $\Lambda=60$, $\kappa_{\max}=30$, $N=91$,
 and the consequence $\eta_Q\le4.9\times10^{-15}$, margin $\le5.9\times10^{-4}$ relative on $\Phi$
 at every tabulated $\zeta$.  This is a genuine, reproducible advance over
 \texttt{tail\_bound.tex}, \S8.
\item \textbf{Quotable as a rigorous bracket, modulo $\widehat D$ and the arithmetic of
 \S\ref{ss:noise}.}  $\Phi^{\rm box}\ge\Phi_W$ (data processing; unconditional) and
 $\Phi^{\rm box}-\Phi_W\le\Phi_{22}+\Phi_{\rm off}$.  The \emph{numbers} $\Phi_{22}$
 ($1.94(-9)$ to $8.38(-6)$) are stable to three digits and may be quoted; the numbers
 $\Phi_{\rm off}$ may not, and therefore neither may the widths $0.64$--$1.19\,\%$ nor the
 $1.7$--$2.4\,\%$.  What can be quoted today is the \emph{stable} a priori bracket
 $\Phi_W\le\Phi^{\rm box}\le\Phi_W(1+0.20)$ (measured $17.8$--$19.9\,\%$), a factor
 $4$--$6$ improvement on \texttt{tail\_bound.tex}'s $82$--$154\,\%$, plus the statement that a
 high-precision rerun of \texttt{cert\_tail.py} is expected to bring this to the $1\,\%$ level.
\item \textbf{Not quotable.}  ``$C_\star=6$'', ``$\Phi-\Phi_W\le0.85\zeta^2\kappa_c^{-2}$'',
 ``verified eigenvalue enclosures'', ``ball arithmetic for the fidelity determinant'', and any
 interval containing the box-cutoff term described as certified.
\end{enumerate}

\end{document}
